\subsubsection{插入机会与资源瓶颈的配合}
后悔插入比较一个待配送货箱的最好插入位置与若干替代位置。可用替代方案少、最好位置又明显便宜的货箱，应优先安排，避免其可行插入机会被其他箱子占用。贪心修复强调当前最小增量，后悔修复强调未来替代成本。两种策略从不同角度组织同一个候选池，不需要人为保证某一策略始终最好。

Pisinger与Ropke用统一框架处理多种车辆路径约束，展示了多类移除和插入策略的组合方式\cite{pisinger2007}。本题将插入代价进一步落实到机型重选和双资源接续：同一箱子的目的地相同，插入不同任务后，受影响的载荷序列、交接偏移与充电长度仍可能不同。几何缓存减少重复地形查询，候选筛选减少完整解码次数，而最终入档保持真实物理与时限核验。

自适应权重分配的是计算机会。Turkeš等对自适应大邻域搜索的元分析表明，适应层的增益与算子选择、实现和比较设置有关\cite{turkes2021}。因此本文列明权重更新与预算，用规定起点下的重复运行描述本实例表现，再用独立下界评价保存方案质量。关键任务移除的针对性来自负载和接续关系，而不来自未经验证的“智能优化”标签。

\subsubsection{一次箱组调整怎样传播到完整日程}
考虑将货箱$b$从任务$r$移入$r'$。首先更新两项箱组，删除不再需要访问的服务区，并尝试目标目的地的允许插入位置。对保留路线和兼容机型，重新检查质量、体积与返航能量；每站交接偏移和返航SOC也随箱数及载荷变化重新计算。超限候选在完整排程前即被排除。

随后将更新后的任务放入相应机型序列，按最早可用飞机和满电电池递推开始时刻。原任务充电时间变化会传到后续用同一电池的任务，继而可能再影响后继飞机链。最后把实际交付事件与原始箱号的期限属性相匹配，检查80箱的唯一配送、硬截止和加权迟到。局部移动一个箱子，由此形成“箱组—物理成本—资源释放—后继开始—交付时限”的完整反馈。

上述组织把计算代价分为物理评价、插入筛选、资源排程与箱级检查四部分。无关机型的已验证属性可以保留，受影响机型的后续事件则必须重放。这既利用局部修改的结构，也避免仅比较变化任务自身的时间而遗漏全局影响。
